\documentclass[11pt]{article}
\usepackage[margin=.72in,top=.58in,bottom=.58in]{geometry}
\usepackage{amsmath,amssymb,amsthm,hyperref}
\renewcommand{\normalsize}{\fontsize{10}{12}\selectfont}
\DeclareMathSizes{10}{10}{8}{6}
\hypersetup{colorlinks=true,linkcolor=blue,urlcolor=blue,citecolor=blue}
\setlength{\emergencystretch}{2em}
\newtheorem{theorem}{Theorem}
\newtheorem{corollary}[theorem]{Corollary}
\newtheorem{lemma}[theorem]{Lemma}
\newcommand{\R}{\mathbb R}
\newcommand{\C}{\mathbb C}
\newcommand{\roots}{\operatorname{Roots}}
\newcommand{\mult}{\operatorname{mult}}
\begin{document}
\begin{center}
{\Large\bfseries Root distribution in the Kadar--Yu Gram sequences}\\[.5em]
Research proof and Lean formalization\\
October 26, 2026
\end{center}
\begin{abstract}
We prove squarefreeness, eventual stabilization of the exterior root count,
and density in $[-2,2]$ for the reduced Gram-polynomial sequences appearing
in Morris--Martin Conjecture~4.5.1. The argument applies to every finite real
symmetric head with one appended free path. Removing the monic gcd of the
first two determinants gives the relevant reduced sequence. A
codimension-one eigenvalue-multiplicity bound proves squarefreeness,
including the boundary index; spectral counts give stabilization and density.
The universal Gram model and its complex coefficient and root-multiset
bridges are verified in Lean. The source application uses the mixed Gram
structure and factorization already established by Morris and Martin.
\end{abstract}

\section{The Gram model and the source conclusion}
Let $G$ be a real symmetric $m\times m$ matrix, $m\ge1$, with a distinguished
coordinate placed last. Form $G_s$ by adjoining a path of $s$ new vertices:
the new diagonal is zero, each adjacent path coupling is $1$, and the first
new vertex is coupled only to the distinguished coordinate, also by $1$.
Every entry of the original head $G$ is retained. Put
\begin{equation}\label{def}
 f_s(X)=\det(XI+G_s),\qquad
 g=\gcd_{\rm monic}(f_0,f_1),\qquad p_s=f_s/g\quad(s\ge0),
 \qquad p_{-1}=Xp_0-p_1.
\end{equation}
All gcds and quotients can be taken over either $\R$ or $\C$; the resulting
monic polynomials agree under coefficient extension. Roots are counted with
multiplicity throughout.

\begin{theorem}\label{general}
Every $p_s$, $s\ge-1$, is monic, real-rooted, and squarefree. There are
integers $S,C\ge0$, with $C\le2m$, such that $p_s$ has exactly $C$ roots
outside $[-2,2]$ whenever $s\ge S$. Moreover,
\[
 [-2,2]\subseteq
 \overline{\bigcup_{s\ge0}\{x\in[-2,2]:p_s(x)=0\}}.
\]
The head $G$ may have arbitrary internal couplings.
\end{theorem}

\begin{corollary}\label{source}
Fix $l\ge-1$ and a partition $\lambda\vdash l+2$. For the reduced sequence
$P_n^{(\lambda)}$ of Morris--Martin \cite[Theorem~3.19]{MM}, every
$P_{l+4+k}^{(\lambda)}$, $k\ge-1$, is squarefree. Its number of roots outside
$[-2,2]$ is eventually constant as $k\to\infty$, and
\[
 \bigcup_{k\ge0}\{\alpha\in[-2,2]:P_{l+4+k}^{(\lambda)}(\alpha)=0\}
 \quad\text{is dense in }[-2,2].
\]
Thus all three assertions of \cite[Conjecture~4.5.1]{MM} hold.
\end{corollary}
The source identification proving the corollary is given in Section~\ref{application}.
The classical spectral and polynomial tools below receive no separate novelty claim.

\section{Reduction and the boundary polynomial}
Write $H_s=-G_s$, so $f_s$ is the genuine characteristic polynomial of $H_s$.
The appended path couplings in $H_s$ are $-1$. Expansion at a terminal
vertex gives
\begin{equation}\label{rec}
 f_{s+2}=Xf_{s+1}-f_s.
\end{equation}
Consequently $g$ divides every $f_s$. The quotients $p_s$ are monic, and
\begin{equation}\label{prec}
 p_{s+2}=Xp_{s+1}-p_s,\qquad
 \gcd(p_s,p_{s+1})=1.
\end{equation}
The initial coprimality follows by removing the actual gcd; the recurrence
propagates it. Since each $p_s$ divides a real symmetric characteristic
polynomial, each splits over $\R$.

Let $H_-$ be $H_0$ with the distinguished coordinate deleted. The first
border identity is
\begin{equation}\label{boundary}
 f_1=Xf_0-\det(XI-H_-),\qquad
 \det(XI-H_-)=g(Xp_0-p_1)=gp_{-1}.
\end{equation}
This proves that $p_{-1}$ is a monic, real-split polynomial as well.
It is coprime to $p_0$ by \eqref{prec}. The $0\times0$ minor has
characteristic polynomial $1$, so the statement includes $m=1$.
No coprimality between $g$ and a quotient is used.

\section{Spectral counts and squarefreeness}
For a real symmetric matrix $A$, let $N_A^+(t)$ and $N_A^-(t)$ count its
eigenvalues strictly above and strictly below $t$, respectively.
\begin{lemma}\label{counts}
If $B$ is a principal submatrix of $A$ obtained by deleting $r$ coordinates,
then, for either sign,
\[
 N_B^\pm(t)\le N_A^\pm(t)\le N_B^\pm(t)+r.
\]
For $r=1$, the multiplicities of any real $t$ as a root of the two
characteristic polynomials differ by at most one.
\end{lemma}
\begin{proof}
For the upper bound, project the span of the eigenvectors of $A$ with
eigenvalues greater than $t$ onto the deleted coordinates together with
the greater-than-$t$ eigenspace of $B$. This projection is injective:
a nonzero vector in its kernel would have quadratic form strictly greater
than $t$ times its squared norm when evaluated with $A$, and at most that
quantity when evaluated with $B$. Thus its dimension is at most
$r+N_B^+(t)$. The corresponding positive subspace of $B$, extended by zeros,
gives the lower bound. Apply the same argument to $-A,-B$ for the other sign.

For $r=1$, put
$\delta_\pm=N_A^\pm(t)-N_B^\pm(t)\in\{0,1\}$.
If $m_A(t),m_B(t)$ denote the root multiplicities, counting the entire
spectrum yields
\[
 m_A(t)-m_B(t)=1-\delta_+-\delta_-\in\{-1,0,1\}.
\]
\end{proof}

Apply the multiplicity bound to $H_s$ and $H_{s+1}$. If $t$ is a root
of $p_s$, it is not a root of $p_{s+1}$ by \eqref{prec}. Hence
\[
 \mult_t(p_s)=\mult_t(f_s)-\mult_t(f_{s+1})\le1,
\]
because the fixed factor $g$ contributes the same multiplicity to both.
Real splitting proves that $p_s$ is squarefree. Applying the same argument
to the principal pair $H_-,H_0$, using \eqref{boundary} and
$\gcd(p_{-1},p_0)=1$, proves squarefreeness at $s=-1$.
This proves the first assertion of Theorem~\ref{general}.

\section{Exterior stabilization and interior density}
The free path on $s$ vertices has the distinct eigenvalues
\begin{equation}\label{grid}
 2\cos\frac{j\pi}{s+1},\qquad 1\le j\le s.
\end{equation}
Changing every path coupling from $1$ to $-1$ gives the same characteristic
polynomial. All eigenvalues in \eqref{grid} lie strictly between $-2$ and $2$.
The path is a principal submatrix of $H_s$ with $m$ deleted head coordinates.
Lemma~\ref{counts} therefore bounds both $N_{H_s}^+(2)$ and
$N_{H_s}^-(-2)$ by $m$. Each is nondecreasing when a new vertex is appended.
They are bounded integer sequences, so both eventually stabilize.
The multiset identity
\begin{equation}\label{multiset}
 \roots(f_s)=\roots(g)+\roots(p_s)
\end{equation}
then subtracts a fixed exterior count and proves the stabilization claim.
The strict thresholds exclude the endpoints, as required for roots outside
the closed interval.

For density, fix $-2<a<b<2$ and choose $a<B<b$. The exact grid
\eqref{grid} supplies arbitrarily many distinct path eigenvalues in $(a,B)$.
For completeness, if $0<u<v<1$, choose $M$ so large that
$(R+1)/(M+1)<v-u$. The $R$ consecutive indices starting at
$\lfloor u(M+1)\rfloor+1$ give grid fractions in $(u,v)$.
Apply this with
$u=\arccos(B/2)/\pi$ and $v=\arccos(a/2)/\pi$.

Choose a path having at least $m+\deg g+1$ such eigenvalues. Let $T_s$ be its
matrix. For a symmetric matrix $A$, let
$I_A(a,B)=N_A^+(a)-N_A^+(B)$, the number of eigenvalues in $(a,B]$.
Lemma~\ref{counts} at both thresholds gives
\begin{equation}\label{interval}
 I_{H_s}(a,B)\ge I_{T_s}(a,B)-m\ge\deg g+1.
\end{equation}
At most $\deg g$ of these roots can be contributed by $g$, counting
multiplicity. By \eqref{multiset}, $p_s$ has a root in $(a,B]\subset(a,b)$.
Every open subinterval of $(-2,2)$ is therefore hit. Passing to closure and
using $\overline{(-2,2)}=[-2,2]$ includes both endpoints. This completes
Theorem~\ref{general}.

\section{Identification of the source sequence}\label{application}
Set $N=l+4$ and $d=\dim S^\lambda\ge1$. Use the real orthonormal Specht
basis in \cite[Section~3.2.3]{MM}. In the mixed family of \cite[(3.17)]{MM},
fix $d-1$ tail parameters at $N$ and let the remaining parameter be $N+s$.
Reorder the attaching vector $u_{l+3,l+4}b_d$ last in the initial matrix.
Equations~(3.2.3), (3.2.4), and the literal border~(3.2.17) show that the
resulting Gram matrices are exactly $XI+G_s$ for one fixed finite real
symmetric head $G$. All interactions between different heads remain in $G$.

Proposition~3.25 of \cite{MM} gives, in this orthonormal convention,
\begin{equation}\label{sourcefactor}
 f_s=C^{(\lambda)}\bigl(P_N^{(\lambda)}\bigr)^{d-1}P_{N+s}^{(\lambda)}
     =F P_{N+s}^{(\lambda)},
\end{equation}
where $F$ is fixed and monic. Reducedness and the minus recurrence imply
$\gcd(P_N^{(\lambda)},P_{N+1}^{(\lambda)})=1$. Thus the monic gcd of
$f_0,f_1$ is exactly $F$, including all of its multiplicities, and
\[
 p_s=P_{N+s}^{(\lambda)}\quad(s\ge0),\qquad
 p_{-1}=P_{N-1}^{(\lambda)}.
\]
The construction agrees with the maximal-common-factor reduction
in \cite[(3.2.12)]{MM}. The source's actual determinant recurrences
(3.2.7) and (3.2.9) have the minus sign used throughout this proof.
Theorem~\ref{general} proves Corollary~\ref{source} for every partition and
all $k\ge-1$.

\section{Formal certificate and scope}
The accompanying Lean development contains 190 declarations in 22 modules.
Its closed endpoints are
\begin{center}
\texttt{gramModelReduced\_root\_distribution}\\
\texttt{gramModelReducedComplex\_root\_distribution}.
\end{center}
They quantify over every
finite real symmetric head and use the literal determinant and normalized
gcd quotients in \eqref{def}. No factorization, consecutive coprimality,
cyclic-compression, or spectral-bound assertion is left as a premise of
these endpoints.

The complex development defines the quotient using the gcd over $\C$,
proves its equality to the real quotient under coefficient mapping,
transfers the entire root multiset, proves every complex root real,
and verifies the exterior count and density with actual complex evaluation.
The boundary polynomial is included. A separate monic-factor uniqueness
lemma identifies any complex source factorization \eqref{sourcefactor}
with this quotient sequence.

The source application in Section~\ref{application} invokes the mixed-Gram
structure and Proposition~3.25 already proved in \cite{MM}; Specht modules
are not reimplemented in the certificate. The formal theorem proves the
universal matrix model that contains those mixed families. Source fidelity
and the scope of this passage are recorded explicitly in the accompanying
review materials.
The strict replay uses Lean~4.34.1 and mathlib revision \texttt{d13f23b7},
compiles every module afresh with warnings as errors, and checks every
declaration's axiom dependencies. Only \texttt{propext},
\texttt{Classical.choice}, and \texttt{Quot.sound} occur.

The proposed source resolution is one combined outcome. Cauchy comparison,
Chebyshev formulas, gcd reduction, the complex scalar extension, and individual
clauses receive no separate credit. The source interpretation and a bounded
prior-art investigation remain subject to independent review; this draft makes
no claim of absolute worldwide priority.

\begin{thebibliography}{9}
\bibitem{MM} B.~Morris and P.~Martin,
\emph{On semisimplicity criteria and non-semisimple representation theory
for the Kadar-Yu algebras}, arXiv:2512.24535v1 (31 December 2025),
\url{https://arxiv.org/abs/2512.24535}.
\end{thebibliography}
\end{document}
